\documentclass[11pt]{article}
\usepackage[margin=1in]{geometry}
\usepackage{amsmath,amssymb,amsthm,mathtools}
\usepackage[T1]{fontenc}
\usepackage[utf8]{inputenc}
\usepackage{enumitem}
\usepackage{booktabs,longtable,array}
\usepackage{hyperref}
\usepackage{xcolor}
\usepackage{microtype}
\hypersetup{colorlinks=true, linkcolor=blue!50!black, citecolor=blue!50!black, urlcolor=blue!50!black}

\newtheorem{theorem}{Theorem}[section]
\newtheorem{proposition}[theorem]{Proposition}
\newtheorem{lemma}[theorem]{Lemma}
\newtheorem{corollary}[theorem]{Corollary}
\newtheorem{claim}[theorem]{Claim}
\theoremstyle{definition}
\newtheorem{definition}[theorem]{Definition}
\newtheorem{assumption}[theorem]{Assumption}
\theoremstyle{remark}
\newtheorem{remark}[theorem]{Remark}

\newcommand{\KL}{\mathrm{KL}}
\newcommand{\Ent}{\mathrm{Ent}}
\newcommand{\Prob}{\mathsf{P}}
\newcommand{\E}{\mathbb{E}}
\newcommand{\R}{\mathbb{R}}
\newcommand{\N}{\mathbb{N}}
\newcommand{\Pcal}{\mathcal{P}}
\newcommand{\Mcal}{\mathcal{M}}
\newcommand{\Qcal}{\mathcal{Q}}
\newcommand{\Zcal}{\mathcal{Z}}
\newcommand{\Fcal}{\mathcal{F}}
\newcommand{\Gcal}{\mathcal{G}}
\newcommand{\W}{\mathcal{W}}
\newcommand{\one}{\mathbf{1}}
\newcommand{\dd}{\,d}
\newcommand{\supp}{\operatorname{supp}}
\newcommand{\argmin}{\operatorname*{arg\,min}}
\newcommand{\argmax}{\operatorname*{arg\,max}}
\newcommand{\esssup}{\operatorname*{ess\,sup}}
\newcommand{\ri}{\operatorname{ri}}
\newcommand{\dom}{\operatorname{dom}}
\newcommand{\Law}{\operatorname{Law}}
\newcommand{\Var}{\operatorname{Var}}
\newcommand{\Cov}{\operatorname{Cov}}
\newcommand{\diag}{\operatorname{diag}}
\newcommand{\Id}{\mathrm{Id}}
\newcommand{\Lip}{\operatorname{Lip}}
\newcommand{\osc}{\operatorname{osc}}
\newcommand{\diam}{\operatorname{diam}}
\newcommand{\proj}{\operatorname{proj}}
\newcommand{\e}{\mathrm{e}}

\newenvironment{argument}{\paragraph{Argument.}\begin{enumerate}[leftmargin=2em]}{\end{enumerate}}
\newenvironment{proofoutline}{\paragraph{Proof architecture.}\begin{enumerate}[leftmargin=2em]}{\end{enumerate}}


% Source-faithfulness apparatus used by the paper reconstructions.
\newcommand{\sourceanchor}[1]{\par\smallskip\noindent\textbf{Source anchor.} #1\par\smallskip}
\newcommand{\sourcecomment}[1]{\begin{remark}#1\end{remark}}
\newenvironment{sourceguide}
  {\begin{description}[style=nextline,leftmargin=3.2cm,labelwidth=3cm]}
  {\end{description}}

% Backward-compatible environments used by some of the older reconstructions.
\newenvironment{distilled}[1][]{\begin{theorem}[#1]}{\end{theorem}}
\newenvironment{distillation}{\begin{enumerate}[leftmargin=2em]}{\end{enumerate}}

\newenvironment{commentarynotes}{\begin{remark}\begin{enumerate}[leftmargin=2em]}{\end{enumerate}\end{remark}}

\title{Equivalence and Singularity of Infinite Product Measures}
\author{}
\date{Kakutani (1948), reconstructed through the supplied modern exposition}
\begin{document}
\maketitle

\section*{Source map and scope}
\begin{description}[style=nextline,leftmargin=3.2cm,labelwidth=3cm]
\item[Primary paper] Shizuo Kakutani, ``On Equivalence of Infinite Product Measures,'' 1948.
\item[Local detailed source] \texttt{\detokenize{Undergrad.2023S - Kakutani measure equivalence.pdf}}
\item[Coverage] Theorem 2.3, Lemmas 2.4--2.8, Corollaries 2.10--2.12, and the Cantor-space Propositions 3.2 and 3.5 from the supplied exposition; these reconstruct Kakutani's product-measure dichotomy and density formula.
\item[Source anchors] equations (2.1), (2.6), (2.9), (2.13)--(2.14), and (3.1)--(3.4).
\end{description}

The supplied source is a modern proof-level exposition rather than a scan of the 1948 article.  The
argument below follows that exposition's order and techniques: finite-product Radon--Nikodym
factorization, square-root density products, an \(L^2\) Cauchy argument, Carath\'eodory extension,
and a small-set proof of singularity.  The final almost-sure infinite-product formula uses the
independent-series result cited by the exposition to van Kampen.

\section{Hellinger affinity}
Let \(\mu,\nu\) be probability measures on \((X,\mathcal B)\), both dominated by \(\lambda\).  Define
\begin{equation}\label{eq:kak-rho}
 \rho(\mu,\nu)=\int_X
 \sqrt{\frac{d\mu}{d\lambda}\frac{d\nu}{d\lambda}}\,d\lambda.
\end{equation}
The Radon--Nikodym chain rule shows that this is independent of \(\lambda\).  If \(\nu\ll\mu\),
\begin{equation}\label{eq:kak-rho-rn}
 \rho(\mu,\nu)=\int_X\sqrt{\frac{d\nu}{d\mu}}\,d\mu.
\end{equation}
Cauchy--Schwarz gives \(0\le\rho\le1\); for equivalent probabilities, \(\rho>0\), and
\(\rho=1\) exactly when \(\mu=\nu\).

The squared Hellinger distance used in the exposition is
\begin{equation}\label{eq:kak-hellinger}
 d_H^2(\mu,\nu)=2(1-\rho(\mu,\nu)).
\end{equation}

\section{Infinite products and the main theorem}
Let \((X_n,\mathcal B_n)\) be measurable spaces and let \(\mu_n\sim\nu_n\) be probability
measures.  Put
\[
 X=\prod_{n\ge1}X_n,
 \qquad
 \mu=\bigotimes_{n\ge1}\mu_n,
 \qquad
 \nu=\bigotimes_{n\ge1}\nu_n.
\]

\begin{theorem}[Kakutani dichotomy]
The following are equivalent:
\begin{enumerate}[leftmargin=2em]
\item \(\mu\) and \(\nu\) are not mutually singular;
\item \(\mu\sim\nu\);
\item \(\prod_{n\ge1}\rho(\mu_n,\nu_n)>0\);
\item \(\sum_{n\ge1}-\log\rho(\mu_n,\nu_n)<\infty\);
\item \(\sum_{n\ge1}d_H^2(\mu_n,\nu_n)<\infty\).
\end{enumerate}
Consequently the two product measures are either equivalent or mutually singular.
\end{theorem}

The source proves the chain
\[
 (3)\Longrightarrow(2)\Longrightarrow(1)\Longrightarrow(3),
 \qquad
 (3)\Longleftrightarrow(4)\Longleftrightarrow(5).
\]
The equivalence of (3) and (4) is the definition of convergence of a positive infinite product.
For numbers \(a_n\in(0,1]\),
\[
 \sum_n(1-a_n)<\infty
 \quad\Longleftrightarrow\quad
 \sum_n-\log a_n<\infty,
\]
because \(-\log a/(1-a)\to1\) as \(a\uparrow1\).  Combining this with
\eqref{eq:kak-hellinger} proves (4)\(\Leftrightarrow\)(5).

\section{Coordinate densities: Lemma 2.4}
Let
\begin{equation}\label{eq:kak-Rn}
 R_n=\frac{d\nu_n}{d\mu_n}.
\end{equation}
Viewed on the full product, \(R_n(x)=R_n(x_n)\) is measurable and the family
\((R_n)\) is independent under \(\mu\), since each variable depends on only one coordinate.
The same observation applies to every finite product.

\section{Finite products: Lemma 2.5}
Write
\[
 \mu_{\le k}=\bigotimes_{n=1}^k\mu_n,
 \qquad
 \nu_{\le k}=\bigotimes_{n=1}^k\nu_n.
\]
For a measurable rectangle \(E=\prod_{n=1}^kE_n\), Tonelli and independence give
\begin{align*}
 \nu_{\le k}(E)
 &=\prod_{n=1}^k\int_{E_n}R_n\,d\mu_n\\
 &=\int_E\prod_{n=1}^kR_n\,d\mu_{\le k}.
\end{align*}
The measure on the right agrees with \(\nu_{\le k}\) on the generating rectangles, hence on the
product sigma-algebra.  Therefore
\begin{equation}\label{eq:kak-finite-rn}
 \frac{d\nu_{\le k}}{d\mu_{\le k}}=\prod_{n=1}^kR_n.
\end{equation}
By symmetry \(\mu_{\le k}\sim\nu_{\le k}\).  Taking square roots and integrating yields
\begin{equation}\label{eq:kak-finite-affinity}
 \rho(\mu_{\le k},\nu_{\le k})
 =\prod_{n=1}^k\rho(\mu_n,\nu_n).
\end{equation}

\section{Square-root densities: Lemma 2.7}
Define
\begin{equation}\label{eq:kak-psi}
 \psi_k=\prod_{n=1}^k\sqrt{R_n}.
\end{equation}
Then \(\psi_k\in L^2(\mu)\), and by Tonelli
\[
 \|\psi_k\|_2^2
 =\int\prod_{n=1}^kR_n\,d\mu
 =\prod_{n=1}^k\int R_n\,d\mu_n=1.
\]
For \(l>k\), independence gives
\begin{align*}
 \langle\psi_k,\psi_l\rangle
 &=\prod_{n=k+1}^l\int\sqrt{R_n}\,d\mu_n
 =\prod_{n=k+1}^l\rho(\mu_n,\nu_n),
\end{align*}
and hence
\begin{equation}\label{eq:kak-cauchy}
 \|\psi_l-\psi_k\|_2^2
 =2\left(1-\prod_{n=k+1}^l\rho(\mu_n,\nu_n)\right).
\end{equation}

\section{Positive affinity product implies equivalence: Lemma 2.8 and Theorem 2.3}
Assume \(\prod_n\rho(\mu_n,\nu_n)>0\).  Tail products then tend to one, so
\eqref{eq:kak-cauchy} makes \((\psi_k)\) Cauchy in \(L^2(\mu)\).  Let
\[
 \psi_k\longrightarrow\psi\qquad\text{in }L^2(\mu).
\]
For a cylinder set \(E\) depending on the first \(m\) coordinates and \(k\ge m\),
\[
 \nu(E)=\int_E\psi_k^2d\mu.
\]
Moreover,
\[
 \|\psi_k^2-\psi^2\|_1
 \le\|\psi_k-\psi\|_2\|\psi_k+\psi\|_2\longrightarrow0.
\]
Thus
\begin{equation}\label{eq:kak-cylinder-limit}
 \nu(E)=\int_E\psi^2d\mu
\end{equation}
for every elementary cylinder.  Define \(\lambda(B)=\int_B\psi^2d\mu\).  By the uniqueness part
of the Carath\'eodory extension theorem, \(\lambda=\nu\) on the whole product sigma-algebra.  Hence
\begin{equation}\label{eq:kak-product-rn}
 \nu\ll\mu,
 \qquad
 \frac{d\nu}{d\mu}=\psi^2.
\end{equation}
Repeating the argument with \(d\mu_n/d\nu_n\) gives \(\mu\ll\nu\), proving equivalence.

\section{Zero affinity product implies singularity}
Assume \(\prod_n\rho(\mu_n,\nu_n)=0\).  Fix \(\varepsilon>0\) and choose \(k\) with
\(\prod_{n=1}^k\rho(\mu_n,\nu_n)<\varepsilon\).  Let
\[
 B_{\le k}=
 \left\{x:\prod_{n=1}^kR_n(x_n)>1\right\}.
\]
On \(B_{\le k}\), \(1\le\sqrt{\prod R_n}\), so
\[
 \mu_{\le k}(B_{\le k})
 \le\rho(\mu_{\le k},\nu_{\le k})<\varepsilon.
\]
On the complement, \(\prod R_n\le\sqrt{\prod R_n}\), whence
\[
 \nu_{\le k}(X_{\le k}\setminus B_{\le k})
 \le\rho(\mu_{\le k},\nu_{\le k})<\varepsilon.
\]
Lift \(B_{\le k}\) to the full product.  For every \(\varepsilon\) there is a measurable \(B\) with
\(\mu(B)<\varepsilon\) and \(\nu(X\setminus B)<\varepsilon\).  Choosing a summable sequence of
errors and taking a limsup/liminf set produces a set of full \(\nu\)-measure and zero
\(\mu\)-measure.  Thus \(\mu\perp\nu\).  This proves the contrapositive of (1)\(\Rightarrow\)(3).

\section{Corollaries 2.10--2.12}
The first corollary is the equivalence/singularity dichotomy itself.  The second is the exact affinity
factorization
\begin{equation}\label{eq:kak-infinite-affinity}
 \rho(\mu,\nu)=\prod_{n\ge1}\rho(\mu_n,\nu_n).
\end{equation}
In the singular case both sides are zero.  In the equivalent case,
\[
 \rho(\mu,\nu)=\int\psi\,d\mu
 =\lim_{k\to\infty}\int\psi_kd\mu
 =\lim_{k\to\infty}\prod_{n=1}^k\rho(\mu_n,\nu_n).
\]

For the pointwise density formula, \(L^2\)-convergence implies convergence in measure.  Since
\(\psi>0\) almost surely in the equivalent case,
\[
 \log\psi_k=\frac12\sum_{n=1}^k\log R_n
\]
converges in measure to \(\log\psi\).  The exposition invokes van Kampen's theorem that a series of
independent random variables converging in measure converges almost surely.  Therefore
\begin{equation}\label{eq:kak-pointwise}
 \psi_k\to\psi\quad\mu\text{-a.s.},
 \qquad
 \frac{d\nu}{d\mu}
 =\lim_{k\to\infty}\prod_{n=1}^k\frac{d\nu_n}{d\mu_n}
 \quad\mu\text{-a.s.}
\end{equation}
The displayed derivative direction is the one consistent with \(R_n=d\nu_n/d\mu_n\).  The supplied
exposition contains a reversed derivative in one displayed line near Corollary 2.12; its subsequent
Cantor-space Proposition 3.5 and the proof identify the intended orientation as
\eqref{eq:kak-pointwise}.

\section{Cantor-space specialization}
Let \(X_n=\{0,1\}\), with
\[
 \mu_n(0)=\alpha_n,
 \qquad
 \nu_n(0)=\beta_n.
\]
Then
\begin{equation}\label{eq:kak-bernoulli-rho}
 \rho(\mu_n,\nu_n)
 =\sqrt{\alpha_n\beta_n}
  +\sqrt{(1-\alpha_n)(1-\beta_n)}.
\end{equation}
The plus sign is forced by the definition of affinity; the supplied PDF displays a minus sign in one
line, a typographical error corrected by the immediately following Hellinger calculation.  Indeed,
\begin{equation}\label{eq:kak-bernoulli-hellinger}
 d_H^2(\mu_n,\nu_n)
 =(\sqrt{\alpha_n}-\sqrt{\beta_n})^2
 +(\sqrt{1-\alpha_n}-\sqrt{1-\beta_n})^2.
\end{equation}

\begin{proposition}[Proposition 3.2]
If some \(\gamma\in(0,1)\) satisfies
\(\gamma\le\alpha_n,\beta_n\le1-\gamma\) for every \(n\), then
\[
 \mu\sim\nu
 \quad\Longleftrightarrow\quad
 \sum_n(\alpha_n-\beta_n)^2<\infty.
\]
Divergence of the sum is equivalent to mutual singularity.
\end{proposition}

\begin{proof}
Rationalize both square-root differences:
\[
 (\sqrt{\alpha_n}-\sqrt{\beta_n})^2
 =\frac{(\alpha_n-\beta_n)^2}
       {(\sqrt{\alpha_n}+\sqrt{\beta_n})^2},
\]
and similarly for the complementary probabilities.  The uniform lower and upper bounds make both
denominators bounded above and below by positive constants.  Thus the Hellinger summand in
\eqref{eq:kak-bernoulli-hellinger} is comparable to \((\alpha_n-\beta_n)^2\), and Kakutani's
criterion applies.
\end{proof}

When \(\mu\sim\nu\), Proposition 3.5 specializes \eqref{eq:kak-pointwise} to
\[
 \frac{d\nu}{d\mu}(x)
 =\lim_{k\to\infty}
 \prod_{n=1}^k\frac{\nu_n(\{x_n\})}{\mu_n(\{x_n\})}.
\]

\section{Gaussian specialization}
For \(\mu_n=N(0,\sigma_n^2)\) and \(\nu_n=N(m_n,\sigma_n^2)\), direct Gaussian integration gives
\[
 \rho(\mu_n,\nu_n)=
 \exp\!\left(-\frac{m_n^2}{8\sigma_n^2}\right).
\]
Hence
\[
 \bigotimes_n\mu_n\sim\bigotimes_n\nu_n
 \quad\Longleftrightarrow\quad
 \sum_n\frac{m_n^2}{\sigma_n^2}<\infty.
\]
This application is the product-coordinate form of the Cameron--Martin square-summability
criterion; the connection is supplemental commentary, not part of the supplied Cantor-space note.

\section{Source-faithfulness remarks}
\begin{enumerate}[leftmargin=2em]
\item The proof of equivalence is an \(L^2\) square-root-density proof, not a direct martingale
convergence proof for \(\prod R_n\).
\item Almost-sure convergence of the infinite product is obtained only after invoking the cited
independent-series theorem of van Kampen.
\item The source's two typographical inconsistencies---the derivative orientation near Corollary
2.12 and the sign in the Bernoulli affinity---are recorded above rather than silently attributed to
Kakutani.
\end{enumerate}

\begin{thebibliography}{99}
\bibitem{ChenGeorgiouPavon2021} Yongxin Chen, Tryphon T. Georgiou, and Michele Pavon. \emph{Stochastic Control Liaisons: Richard Sinkhorn Meets Gaspard Monge on a Schr\"odinger Bridge}. SIAM Review, 63(2):249--313, 2021. doi:10.1137/20M1339982. arXiv:2005.10963.
\bibitem{Leonard2014Survey} Christian L\'eonard. \emph{A Survey of the Schr\"odinger Problem and Some of Its Connections with Optimal Transport}. Discrete and Continuous Dynamical Systems--A, 34(4):1533--1574, 2014. doi:10.3934/dcds.2014.34.1533. arXiv:1308.0215.
\bibitem{Fortet1940} Robert Fortet. \emph{Resolution d'un systeme d'equations de M. Schrodinger}. Journal de mathematiques pures et appliquees, 9e serie, 19(1--4):83--105, 1940.
\bibitem{SinkhornKnopp1967} Richard Sinkhorn and Paul Knopp. \emph{Concerning Nonnegative Matrices and Doubly Stochastic Matrices}. Pacific Journal of Mathematics, 21(2):343--348, 1967.
\bibitem{FranklinLorenz1989} Joel Franklin and Jens Lorenz. \emph{On the Scaling of Multidimensional Matrices}. Linear Algebra and its Applications, 114/115:717--735, 1989. doi:10.1016/0024-3795(89)90490-4.
\bibitem{Carroll2004} Joseph E. Carroll. \emph{Birkhoff's Contraction Coefficient}. Linear Algebra and its Applications, 389:227--234, 2004. doi:10.1016/j.laa.2004.02.039.
\bibitem{EvesonNussbaum1995} Simon P. Eveson and Roger D. Nussbaum. \emph{An Elementary Proof of the Birkhoff--Hopf Theorem}. Mathematical Proceedings of the Cambridge Philosophical Society, 117(1):31--55, 1995.
\bibitem{BlanchetMurthy2019} Jose Blanchet and Karthyek R. A. Murthy. \emph{Quantifying Distributional Model Risk via Optimal Transport}. Mathematics of Operations Research, 44(2):565--600, 2019. doi:10.1287/moor.2018.0936. arXiv:1604.01446.
\bibitem{GaoKleywegt2023} Rui Gao and Anton J. Kleywegt. \emph{Distributionally Robust Stochastic Optimization with Wasserstein Distance}. Mathematics of Operations Research, 48(2):603--655, 2023. doi:10.1287/moor.2022.1275. arXiv:1604.02199.
\bibitem{MohajerinEsfahaniKuhn2018} Peyman Mohajerin Esfahani and Daniel Kuhn. \emph{Data-Driven Distributionally Robust Optimization Using the Wasserstein Metric: Performance Guarantees and Tractable Reformulations}. Mathematical Programming, 171(1--2):115--166, 2018. doi:10.1007/s10107-017-1172-1. arXiv:1505.05116.
\bibitem{WangGaoXie2025} Jie Wang, Rui Gao, and Yao Xie. \emph{Sinkhorn Distributionally Robust Optimization}. Operations Research, 2025. doi:10.1287/opre.2023.0294. arXiv:2109.11926.
\bibitem{Girsanov1960} I. V. Girsanov. \emph{On Transforming a Certain Class of Stochastic Processes by Absolutely Continuous Substitution of Measures}. Theory of Probability and Its Applications, 5(3):285--301, 1960. doi:10.1137/1105027.
\bibitem{CameronMartin1945} R. H. Cameron and W. T. Martin. \emph{Transformations of Wiener Integrals under a General Class of Linear Transformations}. Transactions of the American Mathematical Society, 58(2):184--219, 1945.
\bibitem{Yeh1978} J. Yeh. \emph{Transformation of Conditional Wiener Integrals under Translation and the Cameron--Martin Translation Theorem}. Tohoku Mathematical Journal, 30(4):505--515, 1978.
\bibitem{Kakutani1948} Shizuo Kakutani. \emph{On Equivalence of Infinite Product Measures}. Annals of Mathematics, 49(1):214--224, 1948.
\bibitem{ShiDeBortoliCampbellDoucet2023} Yuyang Shi, Valentin De Bortoli, Andrew Campbell, and Arnaud Doucet. \emph{Diffusion Schr\"odinger Bridge Matching}. Advances in Neural Information Processing Systems, 2023. arXiv:2303.16852.
\bibitem{LiuLipmanNickelKarrerTheodorouChen2024} Guan-Horng Liu, Yaron Lipman, Maximilian Nickel, Brian Karrer, Evangelos A. Theodorou, and Ricky T. Q. Chen. \emph{Generalized Schr\"odinger Bridge Matching}. International Conference on Learning Representations, 2024. arXiv:2310.02233.
\bibitem{DomingoEnrichDrozdzalKarrerChen2025} Carles Domingo-Enrich, Michal Drozdzal, Brian Karrer, and Ricky T. Q. Chen. \emph{Adjoint Matching: Fine-Tuning Flow and Diffusion Generative Models with Memoryless Stochastic Optimal Control}. arXiv preprint arXiv:2409.08861, 2025.

\bibitem{Birkhoff1957} Garrett Birkhoff. \emph{Extensions of Jentzsch's Theorem}. Transactions of the American Mathematical Society, 85(1):219--227, 1957.
\bibitem{Sinkhorn1964} Richard Sinkhorn. \emph{A Relationship Between Arbitrary Positive Matrices and Doubly Stochastic Matrices}. Annals of Mathematical Statistics, 35(2):876--879, 1964.
\bibitem{Sinkhorn1967Prescribed} Richard Sinkhorn. \emph{Diagonal Equivalence to Matrices with Prescribed Row and Column Sums}. American Mathematical Monthly, 74(4):402--405, 1967.
\bibitem{Bauer1965} Friedrich L. Bauer. \emph{An Elementary Proof of the Hopf Inequality for Positive Operators}. Numerische Mathematik, 7:331--337, 1965.
\bibitem{Bushell1973} P. J. Bushell. \emph{Hilbert's Metric and Positive Contraction Mappings in a Banach Space}. Archive for Rational Mechanics and Analysis, 52:330--338, 1973.

\end{thebibliography}

\end{document}
